\scanpage{638}
\mspage{26}

Posons alors

\[
 C=\operatorname{Mono}(M),\quad TC=C\cap W,\quad TF=C_*,\quad F=(TC)_*.\tag*{(2.6.4)}
\]

Sous ces conditions

$1^\circ$) $C\mathrel{\leftarrow\!\!|\!\!\rightarrow}TF$, plus précisément $C\unsure{\mathrel{\boldsymbol{\leftarrow\!\!|\!\!\rightarrow}}|\mathrel{\Leftarrow\!\!|\!\!\Rightarrow}}TF$
est un couple de Quillen spécial.

$2^\circ$) Pour que $TC\longleftrightarrow F$ soit un couple
de Quillen, il f. et s. que $TC$ soit
stable par cochangement de base, et alors
c'est un couple de Quillen spécial.

$3^\circ$) Pour que $(W,C,F)$ soit un triple de
Quillen (lequel sera nécessairement spécial
par $1^\circ,2^\circ$, \deleted{\unsure{donc}} et en particulier \unsure{clos}) il
f. et s. que les conditions suivantes soient
satisfaites

\[
\left\{\begin{aligned}
 \text{a)}\quad & TC\bigl(\overset{\mathrm{def}}{=}W\cap\mathrm{Mono}\bigr)\ \text{est stable par cochange-}\\
 &\quad\text{ment de base}\\
 \text{b)}\quad & TF\bigl(\overset{\mathrm{def}}{=}(\mathrm{Mono})_*
 \unsure{=|\underset{\vert}{=}}\text{ ``applications de}\\
 &\hspace{8em}\text{type ``injectif''}\bigr)\subset W.
\end{aligned}\right.\tag*{(2.6.5)}
\]

% slanted
\margin{\uline{NB} les conditions (2.6.5) équivalent à dire que $(W,C)$ est une structure de \inserted{\unsure{cat. à}} cofibrations sur $M$.}

et la condition b) équivaut encore à \deleted{celle}.

(2.6.6) \uline{b'}) Toute flèche $f$ de $M$ (ou seulement
de $TF$) se factorise en $f=pi$, avec
$i$ mono, $p\in W$. \deleted{\unsure{(On verra en}}

\uline{Corollaire 2.7.} \deleted{\illegible{1}}\inserted{$M$ et $W$ étant comme dans (2.6.1)\deleted{\unsure{a}}, (2.6.3),}
\deleted{\unsure{présence de}} Pour que $(W,C,F)$ soit un triple de
Quillen \uline{propre à gauche} (i.e. $C\subset\operatorname{Cof}_W$), il
f. et s. que $W$ satisfasse la condition
(2.6.5\deleted{\unsure{a}} b) (i.e. \deleted{\unsure{$TF$}} $W$ contient les \unsure{flèches} de
type injectif) et de plus
% archive sheet number: 637
